\chapter{The sequence}
\label{ch:sequence}

\section{Why order is the whole content}

Chapter~\ref{ch:ledgerch} found that almost every measure proposed in this book
has already been tried in Cuba, that most of them worked, and that all of them
were reversed. Part III has added institutional machinery to stop the reversing.
What remains is the order, and the order is not a presentational matter. Three of
the four failures in Part II were failures of sequence rather than of direction.

\emph{1993--95} got the sequence right and is the counter-example: stabilise
first, then liberalise, then grow. It worked, and it was then undone by politics
rather than by economics.

\emph{2011} got it wrong by omission. The Lineamientos legalised activity without
building any of the things that make activity possible --- no legal personality,
no wholesale market, no import rights, no credit, no foreign exchange. The reform
was a permission slip in an economy with no plumbing, and it produced exactly
what such a reform produces.

\emph{2021} got it wrong by inversion. The currency was unified before the
economy was stabilised, before reserves existed, before supply could respond, and
before the private sector was legal. The right measure in the wrong position
destroyed the savings it was meant to rationalise.

So this chapter is the book's actual proposal. Everything before it is inputs.

\section{The four phases}

Dates are relative to a start year rather than to a calendar, because the start
year is not something this book can set.

\begin{figure}[t]
\centering
\begin{tikzpicture}[x=0.275cm,y=1cm,font=\scriptsize]
% 15 years across 8.25cm. The phase shading is a header strip rather than a
% full-height band: banded to the floor it reads as a bar and the rows appear
% to run to year 15, which is the opposite of what the figure says.
\def\rowa{0}\def\rowb{-0.35}\def\rowc{-0.7}
\def\rowd{-1.05}\def\rowe{-1.4}\def\rowf{-1.75}

% phase header
\foreach \a/\b/\n/\g in {0/2/0/8, 2/6/1/16, 6/14/2/8, 14/30/3/16}{
  \fill[black!\g] (\a,0.30) rectangle (\b,0.52);
  \draw[black!45,line width=0.3pt] (\a,0.30) rectangle (\b,0.52);
  \node[font=\tiny\bfseries] at ({(\a+\b)/2},0.41) {\n};}
\node[left,font=\tiny,black!55] at (-0.4,0.41) {phase};

% phase boundaries carried down through the rows
\foreach \x in {2,6,14}{
  \draw[black!25,line width=0.3pt,dotted] (\x,0.28) -- (\x,-1.95);}

% axis
\draw[black!55,line width=0.5pt] (0,-1.95) -- (30,-1.95);
\foreach \x/\l in {0/0,2/1,6/3,14/7,30/15}{
  \draw[black!55,line width=0.5pt] (\x,-1.95) -- (\x,-2.05);
  \node[below,black!55,font=\tiny] at (\x,-2.06) {\l};}
\node[below,black!55,font=\tiny] at (15,-2.34) {years from the start};

\newcommand{\wsbar}[4]{%
  \fill[black!#4] (#2,#1-0.095) rectangle (#3,#1+0.095);
  \draw[black!70,line width=0.4pt] (#2,#1-0.095) rectangle (#3,#1+0.095);}

\wsbar{\rowa}{0}{2}{45}    \wsbar{\rowa}{0.6}{4}{20}  \wsbar{\rowa}{4}{20}{65}
\wsbar{\rowb}{0}{1}{45}    \wsbar{\rowb}{2}{3}{65}    \wsbar{\rowb}{3}{10}{20}
\wsbar{\rowc}{0}{1}{65}    \wsbar{\rowc}{2}{5}{45}    \wsbar{\rowc}{5}{6}{20}
\wsbar{\rowd}{1}{7}{45}    \wsbar{\rowd}{7}{8}{20}
\wsbar{\rowe}{6}{14}{45}
\wsbar{\rowf}{2}{8}{20}    \wsbar{\rowf}{6}{14}{45}   \wsbar{\rowf}{14}{30}{65}

\foreach \y/\l in {\rowa/Energy,\rowb/Money,\rowc/Law,\rowd/Property,
                   \rowe/Welfare,\rowf/Sales}{
  \node[left,font=\tiny] at (-0.4,\y) {\l};}

% the gate: labelled above the rows, where nothing else sits
\draw[black!75,dashed,line width=0.5pt] (10,0.28) -- (10,-1.95);
\node[above right,font=\tiny\itshape,black!75,inner sep=1pt] at (10,0.55)
  {national elections};
\end{tikzpicture}
\caption{The sequence. The header strip marks the four phases; each row is a
workstream and each bar a distinct activity within it, darker shading indicating
the later activity in that row --- not budget, not effort. Energy runs repeals
and efficiency, then critical loads, then the build-out; Money publishes the
debt stock, converts the currency, then restructures; Law repeals and amnesties,
then builds courts, then holds municipal elections; Property registers before it
sells; Welfare converges only in Phase~2; Sales go small, then medium, then ---
after year seven --- large. The ordering claims are the content: the build-out
cannot start before the regulator and the tariff exist, the currency is not
converted before the debt stock is published, the courts precede the first
competitive elections, registration precedes any sale, welfare convergence waits
on revenue that Phase~1's tax administration will not produce until Phase~2, and
the largest assets are sold last --- which is the single clause standing between
this design and the outcome described in Chapter~\ref{ch:integrity}.}
\label{fig:sequence}
\end{figure}

\subsection*{Phase 0 --- months 0 to 12: stop the bleeding}

The organising principle is that Phase 0 contains \emph{only} things that are
fast, cheap, and reversible in their costs but not in their benefits. Nothing
here requires foreign money, a constitutional amendment, or anyone's permission.

The repeals first, because they are free. The right of any resident to contract
with a foreign client and be paid into their own account
(Design~\ref{d:contract}). Rooftop generation by right and duty-free import of
energy and computing equipment (Designs~\ref{d:genrights}.1, \ref{d:genrights}.6,
\ref{d:contract}.5). Activity permitted by right rather than by enumerated list.
Professionals allowed to practise privately. The remaining exit restrictions
lifted. The amnesty of Design~\ref{d:amnesty}.6 --- everyone convicted for
political expression, association or the events of 11 July 2021 released and
their convictions vacated --- which costs the state nothing and is the single
clearest signal available that this time is different. Ra\'ul Castro's most
popular measures were all of this kind and the lesson is in
Chapter~\ref{ch:raul}: begin with the prohibitions that cost nothing to remove,
because they are immediate, they are free, and they buy the credibility the
expensive measures will need.

Then the emergencies. Fuel, spares and maintenance for the existing thermal
fleet, funded as a first call on the budget --- unglamorous, unavoidable, and the
thing a solar enthusiast would skip. Demand-side efficiency at national scale, on
the 2006 precedent. Photovoltaics and storage on critical loads: hospitals, water
pumping, cold stores, bakeries, telecoms. Food imports maintained. Current
supplier invoices paid and the 2009 frozen balances released
(Design~\ref{d:debt}.6).

And the announcements that start the clocks: the complete debt and arrears stock
published; the dollarisation date and conversion rate announced; the political
calendar of Design~\ref{d:calendar} entrenched; the cadastre programme begun.

\subsection*{Phase 1 --- years 1 to 3: build the plumbing}

Money: convert to the dollar at a single rate, all balances at once, wages and
pensions protected at the floor (Design~\ref{d:regime}). Legislate the
prohibition on monetary financing and the recurrent balance rule. Open the debt
restructuring with the claims inside it.

Law: judicial reform and, above all, the administrative court in which a citizen
or a firm can sue the state (Design~\ref{d:courts}.4). Press, association and the
right to information. The party law and the electoral commission. Competitive
municipal elections at the end of the phase.

Property: cadastre and registry delivered; usufruct converted to real,
mortgageable title (Design~\ref{d:tenure}); acopio abolished; wholesale markets
opened; the residential rental market legalised.

Investment: the single investment law with direct employment, deemed consent,
national treatment and arbitration (Design~\ref{d:invest}). Small-scale
privatisation begins --- thousands of retail, catering, repair and service units,
by open auction, worker and cooperative preference.

Energy: the regulator, the standard power purchase agreement, net billing, and
the first competitively tendered independent producer capacity.

And the long-lead item that yields nothing in this phase and must start in it
anyway: the revenue administration. Registration, taxpayer numbers, withholding,
the presumptive regime for micro-enterprise, and the beginning of the VAT build.
Every country that has done this has taken five to eight years, and the welfare
commitment of Chapter~\ref{ch:welfare} is unfinanceable if it slips.

\subsection*{Phase 2 --- years 3 to 7: the growth phase}

The debt restructuring closes and the claims are settled in paper within the cap.
Multilateral membership is pursued. VAT comes into force and the tax base begins
to exist.

The energy build-out reaches scale --- the bulk of the six gigawatts and the
storage --- weighted to the eastern provinces. Tourism reorients on the
value-per-visitor metric, with the moratorium on state hotel construction in
force and the heritage authorities established. The second track becomes
visible: connectivity redundancy delivered, the residence regime open,
service exports counted in the balance of payments. The medicines regulator
reaches stringent-authority standard and the first product partnerships are
signed.

The welfare convergence begins here rather than earlier, because this is the
first phase in which there is revenue: public pay benchmarked and raised under
Design~\ref{d:pay}, the universal cash benefits paid reliably for two full years
before any ration item is withdrawn, and the pension reform legislated with its
indexation formula.

Medium-scale privatisation proceeds under published auction. National assembly
elections fall in this phase, at year five.

\subsection*{Phase 3 --- years 7 to 15: the hard remainder}

What was deliberately left: the large state sector, the utilities, the ports, the
telecommunications incumbent, and the military conglomerate's core assets, sold
only once the competition authority, the sector regulators, the courts, the press
and the beneficial-ownership register are all functioning
(Design~\ref{d:privat}.2, Design~\ref{d:gaesa}.5). Unbundling of the electricity
network. Convergence of benefit levels toward the European standard on the
published trajectory. And, if the conditions of Design~\ref{d:regime}.5 are met,
the question of whether to reintroduce a national currency.

\section{Who carries which phase}

The political-economy question, which a design document is not entitled to skip.

\emph{Phase 0 is executable by the present government}, and this is deliberate.
Every measure in it is within the existing constitutional order, none requires
surrendering the party's leading role, several would relieve pressures the
government is currently absorbing, and the amnesty and the repeals would be
enormously popular at almost no fiscal cost. A book that requires a rupture
before its first measure is a book that does nothing until the rupture, and
Cubans have been waiting for the rupture since 1991.

\emph{Phase 1 requires the party to accept an independent court}, which is the
hardest single concession in the book, because it is the concession from which
the others follow. The coalition that could carry it is the one that already
exists inside the Cuban state: the technocrats who have been writing these
reforms since 2011 and watching them be reversed; the enterprise managers who
cannot operate without a currency that means something; and --- this is the
uncomfortable part --- the armed forces, whose commercial interests are better
served by a functioning economy with enforceable contracts than by an
increasingly worthless monopoly over an increasingly empty one.

\emph{Phase 2 requires the bargain of Design~\ref{d:gaesa}.6} to have been struck
explicitly: security, pensions, property and a legal route into ordinary
political and business life, in exchange for the surrender of discretionary
control. This is where a transition is either agreed or contested, and the book's
position is that it should be negotiated openly and written down, because a
bargain left implicit is renegotiated at every step.

\emph{Phase 3 requires a government formed after the elections of Phase 2.} It is
not the same government, and the design does not assume that it is friendly to
what came before. That is the test of whether the commitment devices of
Chapter~\ref{ch:polity} were real.

\section{What it costs and where it comes from}

Orders of magnitude over a decade, with assumptions stated. Every figure is
marked unverified in the manuscript and listed in \texttt{verify.md}.

\begin{center}
\footnotesize
\begin{tabular}{@{}L{2.3in}r@{}}
\toprule
\multicolumn{2}{@{}l}{\emph{Capital, cumulative over ten years}}\\
Photovoltaics and storage (Ch.~\ref{ch:energy}) & \$6--8~bn \\
Grid transmission and distribution & \$2--3~bn \\
Water: leak reduction and treatment & \$1--2~bn \\
Cadastre, registry, revenue administration & \$0.3--0.5~bn \\
Strategic food reserve, initial & \$0.3--0.5~bn \\
\midrule
\multicolumn{2}{@{}l}{\emph{Recurrent, annual at maturity}}\\
Public pay settlement, health and education first & \$0.6--1.1~bn/yr \\
Benefit convergence, phased & rising \\
\midrule
\multicolumn{2}{@{}l}{\emph{Sources}}\\
Hotel construction redirected & \$1--1.5~bn/yr \\
Avoided fuel imports, from year six & \$1.2--1.5~bn/yr \\
Diaspora investment and remittance conversion & --- \\
Concessional and bilateral finance & --- \\
Privatisation proceeds (to the fund) & --- \\
Tax base as it grows & --- \\
\bottomrule
\end{tabular}
\end{center}

\noindent One line in that table carries most of the argument, and it is best
stated as a share rather than in dollars, because the peso conversion is the
problem the note on sources described. In 2024 Cuba put 37.4 per cent of its
total national investment into tourism-related construction and 2.7 per cent
into agriculture, while running its hotels at occupancy around a quarter,
importing most of its food and appealing for children's milk. Moving even half
of the tourism share into generation, network repair, water and agriculture ---
which requires no foreign money, no legislation and nobody's permission ---
is of the same order as the entire energy programme in this book. The energy
programme then saves, in perpetuity, more per year in foreign exchange than the
hotels were ever going to earn.

\emph{The largest available source of finance for this plan is not foreign. It
is Cuba's own investment budget, spent differently.} That conclusion is the
reason
Chapter~\ref{ch:integrity} matters more than any chapter about capital markets:
the constraint is not the absence of money but the absence of anyone able to
say no to the people currently allocating it.

\section{What could go wrong}

Six failure modes, in descending order of likelihood.

\emph{Capture.} The incumbents run Phases 0 and 1, acquire the assets in Phase 2
before the institutions bite, and Cuba arrives in 2045 as a middle-income country
with a captured state, a compliant press and an ownership structure fixed for a
generation. This is the most likely failure and it is the one that is
irreversible. Everything in Chapter~\ref{ch:integrity} is aimed at it, and the
single indicator that detects it early is Design~\ref{d:measure}.1.

\emph{The workforce leaves anyway.} The reforms take four years to raise incomes
and the exodus takes two to remove the people who would have benefited. Cuba
ends the decade with better institutions and no one to use them. This is why net
migration is the master indicator below and why Phase 0 is judged on whether the
outflow slows rather than on growth.

\emph{Collapse before Phase 1 completes.} A grid failure that is not recovered, a
fuel supplier who stops, a hurricane in the wrong year, a fiscal crisis during
the currency conversion. The design's answer is Phase 0's unglamorous
insistence on keeping the existing fleet running and the food coming while the
new things are built.

\emph{Sequencing failure.} Elections before courts; privatisation before the
registers; the ration withdrawn before the cash arrives; the currency converted
before stabilisation. Each of these has a named clause against it and each will
be argued for, by serious people, at the time.

\emph{The J-curve, politically.} Reforms of this kind hurt before they help.
Prices rise before supply responds; state employment falls before private
employment absorbs it. The political system has to survive the trough, and a
government that faces its first competitive election at the bottom of it may not.
This is an argument for the pacing above --- welfare convergence and elections
both in Phase 2, after the first benefits are visible --- and it is an honest
weakness of any gradualist design.

\emph{An external swing.} A sudden American normalisation would flood a country
with no absorptive capacity, no registry, no planning system and no claims
settlement, and the assets would be allocated in eighteen months to whoever was
positioned. A sharp tightening removes remittances and travel, which are the
household economy. The design's premise is to assume neither, and the practical
protection against the first is simply to have done Phase 1 before it happens.

\section{How to tell whether this is working}

A proposal that cannot be falsified is not a proposal. These are published
annually, and each has a direction that constitutes failure.

\begin{design}{The indicators}
\label{d:indicators}
\begin{rules}
\rl{1}{Net migration} The master indicator. If Cubans continue to leave at
anything like the rate of 2021--25, nothing else in this book matters, because
the design is labour-constrained and the labour is going.

\rl{2}{Load shedding} Average daily hours of interruption, national and by
province, with the eastern provinces reported separately. Everything is
downstream of this.

\rl{3}{Health and education outcomes} Infant mortality, life expectancy, school
attainment, and resident physicians, nurses and teachers per capita ---
\emph{resident}, not counting missions abroad. These must not fall. If they fall,
the design has failed on its own stated terms regardless of any growth figure.

\rl{4}{Real wages and the price of the basket} Not the official index --- the
cost of an actual consumption basket against actual median income.

\rl{5}{Food} Share of consumption domestically produced, and post-harvest losses.

\rl{6}{Service exports excluding medical missions} The second track's existence
test. If this is not visible by year five, Chapter~\ref{ch:talent} was wrong.

\rl{7}{Concentration} Per Design~\ref{d:measure}.1: the share of privatised
assets acquired by the largest ten beneficial owners and by former officeholders.
The early-warning indicator for capture.

\rl{8}{Openness of contracting} Share of public contracts by value awarded
through open tender with more than two bidders.

\rl{9}{Time to authorise} Median days to obtain a business authorisation, and
the share decided within the statutory period.

\rl{10}{Tax revenue} Revenue as a share of output, by instrument. The financing
test for the welfare commitment.

\rl{11}{Reef and protected-area condition} Published annually, per
Design~\ref{d:protect}.6.

\rl{12}{Disaggregation} Every one of the above reported by province and, where
individual, by race. Cuba declared the racial question solved in 1962, made it
unmeasurable for thirty years, and entered its market opening with no idea where
it stood. That is not repeated.
\end{rules}
\end{design}

\section{What this book does not claim}

It does not claim to know what will happen. It is a design, not a forecast, and
the difference is that a design states its assumptions in a form that lets
someone else change them.

It does not claim that the measures here are the only ones that would work.
Several of its recommendations --- dollarisation most obviously, the ordering of
courts before elections next --- are genuinely contestable, and the book has
tried to give the argument on the other side each time rather than to win by
omission.

It does not claim that the embargo is unimportant, only that it is the one
variable Cubans do not control, and that a design which waits on it is not a
design.

It does not claim that the record of 1959--2026 is a simple story. Parts I and II
were written at length precisely because it is not. A state that eliminated
illiteracy in a year also ran labour camps; a state that held infant mortality
through a thirty-five per cent contraction also handed thirty-year sentences
to teenagers for walking in the street. Both libraries in the preface were
telling the truth and neither was being honest, and any proposal that needs one
of them to be false is a proposal that will fail on contact with the country.

And it does not claim that any of this is likely. The most probable future for
Cuba is the continuation of the present: a slow contraction, an ageing remainder,
the lights going out province by province, and the young leaving until there is
nobody left to design for. That future requires no decisions at all, which is
what makes it the default.

What the book claims is narrower and, I think, defensible. Cuba has assets that
compose --- an educated population, a medical and scientific capability with
products and a reputation, an intact environment in a region that has spent its
own, sunlight falling on a country that burns imported oil, and two million of
its own people one flight away with capital and a reason to care. Those assets
have been available for twenty years and have not converted, and the reason is
not principally the embargo. It is that Cuba has never built the machinery that
turns an asset into an investment: a currency that means something, a court that
can be lost in fairly, a registry that says who owns what, a customs service that
does not take a cut, a procurement system that publishes what it bought, and a
rule that cannot be withdrawn the moment it starts to work.

That machinery is buildable. It is mostly not expensive. Almost none of it
requires anyone's permission. And it is the difference between a country with
sunshine and doctors and beaches that is poor, and one that is not.
